\begin{abstract}
Residual networks remain a central architecture family for deep visual recognition, but the
effect of residual-block modifications is best interpreted together with the training recipe
under which they are evaluated.
We study a controlled question on CIFAR-10 at depth~20 with shortcut Option~A: under a frozen
post-activation ResNet-20-A training recipe, does replacing the residual block by a full
pre-activation unit reduce best test error?
We compare Plain-20, ResNet-20-A, PreAct-ResNet-20-A, and targeted BN+ReLU ablations~A/B in a
single-seed, fixed-protocol architecture comparison.
ResNet-20-A improves over Plain-20 ($\ErrResNet\%$ vs.\ $\ErrPlain\%$ best test error;
$\ResidualAdvantagePP$\,pp).
Full PreAct-ResNet-20-A reaches $\ErrPreAct\%$
($\DeltaPreActMinusResNet$\,pp higher error than ResNet-20-A);
\textbf{the hypothesized improvement was not supported under the frozen training recipe and
single-seed setting.}
Ablation~A ($\ErrAblationA\%$) improves on complete PreAct by $\DeltaAblationAVsPreAct$\,pp in
magnitude, but differs from ResNet-20-A by only $\DeltaAblationAMinusResNet$\,pp---a small
single-run gap that we do not treat as a validated improvement.
Ablation~B fails with numerical instability (\AblationBStatus) and is not used as an accuracy
baseline.
These findings do not claim that pre-activation residual units are generally ineffective; they
provide an artifact-backed, protocol-controlled negative result that narrows the hypothesis space
and motivates evaluating architecture changes jointly with optimization settings.
\end{abstract}
